\documentclass[10pt,a4paper]{amsart}
\usepackage[margin=19mm]{geometry}
\usepackage{amsmath,amssymb,mathtools}
\usepackage[scr=rsfs,cal=euler]{mathalfa}
\usepackage{xfrac}
\usepackage[unicode,hidelinks,bookmarks=false]{hyperref}
\newcommand{\ie}{i.e.\ }
\newcommand{\cf}{cf.\ }
\newcommand{\resp}{respectively\ }
\newcommand{\Subsection}{\subsection}
\providecommand{\nobreakdash}{\nobreak\hbox{-}}
\setlength{\emergencystretch}{2em}
\multlinegap0pt
\usepackage{bm}
\usepackage{bbm}
\usepackage{dsfont}
\usepackage{xspace}
\usepackage{cancel}
\usepackage{mathtools}
\usepackage{amsthm}
\makeatletter
\@ifundefined{lem}{\newtheorem{lem}{Lemma}}{}
\@ifundefined{prop}{\newtheorem{prop}{Proposition}}{}
\makeatother
\title[On polynomial automorphisms of Nagata type]{On polynomial automorphisms of Nagata type}
\author{Jorge A. C. Huarcaya, Joe Palacios}
\date{Source: arXiv:2401.12523v2 (CC BY 4.0)}
\begin{document}
\maketitle

\noindent\emph{Source and licence.} On polynomial automorphisms of Nagata type. Version arXiv:2401.12523v2 (CC BY 4.0), distributed under the Creative Commons Attribution 4.0 International licence, as recorded in the arXiv licence metadata for this exact version and shown on the abstract page https://arxiv.org/abs/2401.12523v2. Full rights notice: licenses/arxiv-2401.12523-CC-BY-4.0.txt.\par\smallskip

\noindent\textit{Editorial scope.} Counted unit: the Prop whose source label is \texttt{pr1} in the article (source0.tex lines 449--465, source label \texttt{pr1}), delivered together with the article's own complete original proof of it (source0.tex lines 468--568, one \texttt{proof} environment of 101 lines). No mathematics is written, repaired or re-derived: statement and proof are the article's own text line for line. Required context retained uncounted: eq1 (lines 369--371); psedp (lines 418--422). Every definition, display and label of those lines is retained unchanged. Omitted from this package and not redistributed: 1--368, 372--417, 423--448, 466--467, 569--698. Nothing inside a retained excerpt is omitted or altered: every retained body is delivered exactly as the article's lines are written, so no omission receipt of any kind accompanies this package. Citations: the retained excerpts carry no citation command at all, so no bibliography entry of the article is transcribed and no reference is redistributed. Reference adaptation: none was needed and none was made; every reference inside the delivered unit resolves inside it, so no unresolved reference or citation is produced. Printed slips kept literal: no printed word, symbol, hypothesis or number of the counted statement or of its proof was altered. Layout and class: the article's own class is \texttt{amsart} with packages including inputenc, xy, amsmath, amssymb, amsthm, multicol, amscd, graphics, graphicx, amsaddr, hyperref; the wrapper is the lane's pinned \texttt{amsart} editorial wrapper at 10pt on A4, which restates the theorem-like environments the retained text needs and adds the article's own macro definitions that the retained text needs (none), with extra wrapper packages (bm, bbm, dsfont, xspace, cancel, mathtools). The delivered numbering is the unit's own. Assets: no retained span contains a figure environment, a graphics-inclusion command, a PDF-page inclusion, a table or a TikZ picture; no image file is copied into this package, the package declares no asset dependency and no image file is redistributed with it. Licence: version arXiv:2401.12523v2 of this article is distributed under the Creative Commons Attribution 4.0 International licence, as recorded in the arXiv licence metadata for this exact version and shown on the abstract page https://arxiv.org/abs/2401.12523v2. A full rights notice is delivered at licenses/arxiv-2401.12523-CC-BY-4.0.txt. Characters: every retained excerpt is pure ASCII, so the delivered engine, XeLaTeX with its default Latin Modern text font, reports no missing character.
	\begin{equation}\label{eq1}
		-2y\varphi_x+z\varphi_y=0
	\end{equation}

\begin{lem}\label{psedp}
	Let $\varphi$ be a homogeneous polynomial in $F[x,y,z]$ of degree $n$ that is a solution of the linear PDE \eqref{eq1}. Then, $\phi^{n}=ay^{n}$ and $\phi^{1}=bx$, where $a$ and $b$ are constants in $F$, and moreover
	$$2y\phi^{n-k-1}_x=\phi^{n-k}_y$$
	for all $k=0,1,\dots, n-2$.
\end{lem}

\begin{prop}\label{pr1}
	Let $\varphi\in F[x,y,z]$ be a homogeneous solution of the linear PDE \eqref{eq1}. 
	We have the following assertions:
	\begin{enumerate}
		\item[ ($a$)] If $\varphi$ has degree $2n$, then it has the form: 
		\begin{equation*}
			\varphi=a(xz+y^2)^n+z\psi\,,
		\end{equation*}
		where $a\in F$ and $\psi\in F[x,y,z]$ is a homogeneous polynomial which is also a solution of the equation \eqref{eq1} of degree $2n-1$.
		\item[($b$)] If $\varphi$ has degree $2n+1$,  then it has the form: 
		\begin{equation*}
			\varphi=z\psi\,,
		\end{equation*}
		where $\psi\in F[x,y,z]$ is homogeneous of degree $2n$ which is also a solution of \eqref{eq1} of degree $2n$.
	\end{enumerate}
	
\end{prop}

\begin{proof}
	($a$). Let us suppose that:
	\begin{equation}\label{eqh}
		\varphi=\phi^{2n}+\phi^{2n-1}z+\cdots+\phi^{1}z^{2n-1}+\phi^0z^{2n}\,,
	\end{equation}
	where $\phi^{2n-k}\in F[x,y]$ is a homogeneous polynomial of degree $2n-k$ for  $0\leq k\leq 2n$. By Lemma \ref{psedp}, we get the following equalities: 
	\begin{align}	
		\phi^{2n}&=ay^{2n}\label{eq7}\,\,\,\, \textnormal{and}\,\,\,\, \phi^{1}=bx\\
		2y\phi^{2n-k-1}_x&=\phi^{2n-k}_y\,\,\, \textnormal{para todo}\,\,\, k=0,1,\dots, 2n-2\,.\label{eq8}
	\end{align}
	We claim that there are constants $a, a^{2n-k}_{k-1-j}$, with  $1\leq k\leq n$ y $0\leq j\leq k-1$, such that:
	\begin{align}
		\phi^{2n-k}&=a\binom{n}{k}y^{2n-2k}x^k+\sum_{j=0}^{k-1} a^{2n-k}_{k-1-j}y^{2n-2k+1+j}x^{k-1-j}\label{eq12}
		%	&=a\binom{n}{k}y^{2n-2k}x^k+a^{2n-k}_{k-1}y^{2n-2k+1}x^{k-1}
	\end{align}
	Indeed, taking $k=0$, from the equalities (\ref{eq7}) and  (\ref{eq8}), we find that:
	\begin{equation*}
		2y\phi^{2n-1}_x=\phi^{2n}_y=2nay^{2n-1}\,,
	\end{equation*}
	so that $\phi_x^{2n-1}=nay^{2n-2}$. Hence, integrating with respect to $x$, we obtain
	$$\phi^{2n-1}=nay^{2n-2}x+a^{2n-1}_0y^{2n-1}\,,$$ which can be written as 
	\begin{align*}
		\phi^{2n-1}=a\binom{n}{1}y^{2n-2}x+a^{2n-1}_{0}y^{2n-1}\,,
	\end{align*}
	thus \eqref{eq12} is true for $k=1$. Now, suppose that we have \eqref{eq12}. 
	Taking the derivative with respect to $y$, we have
	$$\phi^{2n-k}_y=2a\binom{n}{k}(n-k)y^{2n-2k-1}x^k+\sum_{j=0}^{k-1}\, (2n-2k+1+j)a^{2n-k}_{k-1-j}y^{2n-2k+j}x^{k-1-j}$$
	By the equality (\ref{eq8}), we get 
	$$\phi^{2n-k-1}_x=a\binom{n}{k}(n-k)y^{2n-2k-2}x^k+\sum_{j=0}^{k-1} \frac{(2n-2k+1+j)a^{2n-k}_{k-1-j}}{2}y^{2n-2k-1+j}x^{k-1-j}\,.$$
	Integrating with respect to $x$, there exists a constant $a^{2n-k-1}_{0}\in F$ such that
	$$\phi^{2n-k-1}=a\binom{n}{k+1}y^{2n-2k-2}x^{k+1}+\sum_{j=0}^{k-1} a^{2n-k-1}_{k-j}y^{2n-2k-1+j}x^{k-j}+a^{2n-k-1}_{0}y^{2n-k-1}\,,$$
	where $\displaystyle a^{2n-k-1}_{k-j}= \frac{(2n-2k+1+j)a^{2n-k}_{k-1-j}}{2(k-j)}$ for all $0\leq j\leq k-1$. Thus, we obtain
	$$\phi^{2n-k-1}=a\binom{n}{k+1}y^{2n-2k-2}x^{k+1}+\sum_{j=0}^{k}a^{2n-k-1}_{k-j}y^{2n-2k-1+j}x^{k-j}\,,$$
	as wished, and this completes the recurrence.
	
	Now, replacing the expressions (\ref{eq7}) and (\ref{eq12}) in the equation (\ref{eqh}), we have  
	\begin{align*}
		\varphi
		&=\phi^{2n}+\sum_{k=1}^{n}\phi^{2n-k}z^k+\sum_{k=n+1}^{2n}\phi^{2n-k}z^k\nonumber\\
		&=ay^{2n}+ \sum_{k=1}^n \left(a\binom{n}{k}y^{2n-2k}x^k+\sum_{j=0}^{k-1} a^{2n-k}_{k-1-j}y^{2n-2k+1+j}x^{k-1-j}\right)z^k +\sum_{k=n+1}^{2n}\phi^{2n-k}z^k\\
		&=a\sum_{k=0}^{n}\binom{n}{k}y^{2n-2k}x^kz^k+ \sum_{k=1}^{n}\left(\sum_{j=0}^{k-1} a^{2n-k}_{k-1-j}y^{2n-2k+1+j}x^{k-1-j}\right)z^k+\sum_{k=n+1}^{2n}\phi^{2n-k}z^k\,.
	\end{align*}
	From  the last equality, we can write
	\begin{equation*}
		\varphi=a(xz+y^2)^n+z\psi, 
	\end{equation*}
	where $\psi\in F[x,y,z]$ is a homogeneous polynomial of degree equal to $2n-1$. By the linearity of the solutions of the equation (\ref{eq1}), we deduce that $z\psi$ is a solution, hence $\psi$ is also a solution.
	
	($b$) Let us suppose that: 
	\begin{equation}\label{eqhi1}
		\varphi=\phi^{2n+1}+\phi^{2n}z+\cdots+\phi^{1}z^{2n}+\phi^0z^{2n+1}\,,
	\end{equation}
	where $\phi^{2n+1-k}$ is a homogeneous  polynomial in $ F[x,y]$ of  degree equal to $2n+1-k$, for all $0\leq k\leq 2n+1$.  By Lemma \ref{psedp}, we get the following equalities:  
	\begin{align}	
		\phi^{2n+1}&=ay^{2n+1}\label{eqi7}\\	
		%\varphi^{1}&=bx\nonumber\\
		2y\phi^{2n-k}_x&=\phi^{2n+1-k}_y\,\,\, \textnormal{para todo}\,\,\, k=0,1,\dots, 2n-1\,.\label{eqi8}
	\end{align}
	As in the even case, we claim that there are constants $a,b_{k-j}^{2n-k}$,  with $0\leq k\leq n-1$ and $0\leq j\leq k$, such that:
	\begin{equation}\label{eqodd}
		\phi^{2n-k}=P^{2n+1}_{k}ay^{2n-2k-1}x^{k+1}+\sum_{j=0}^{k}b^{2n-k}_{k-j}y^{2n-2k+j}x^{k-j}\,,
	\end{equation}
	where 
	$
	P^{2n+1}_{k}=\frac{(2n+1)\cdot(2n+1-2)\cdots (2n+1-2k)}{2^{k+1}\cdot (k+1)!}
	$
	for all index $0\leq k\leq n-1$. 
	
	Replacing $k=0$ in the equation (\ref{eqi7}) and joining with  (\ref{eqi8}), we have
	\begin{equation*} 
		2y\phi^{2n}_x=\phi^{2n+1}_y=(2n+1)ay^{2n}\,,
	\end{equation*} 
	hence $\phi_x^{2n}=\cfrac{(2n+1)}{2}\, \, ay^{2n-1}$.
	Integrating with respect to $x$, we obtain
	\begin{equation*}
		\phi^{2n}=\cfrac{(2n+1)}{2}\,a y^{2n-1}x+b^{2n}_{0}y^{2n}
	\end{equation*}
	thus \eqref{eqodd} is true for $k=0$. 
	Suppose that we have \eqref{eqodd}. Taking the derivative with respect to $y$, we have
	$$\phi^{2n-k}_y=P^{2n+1}_{k}(2n-2k-1)ay^{2n-2k-2}x^{k+1}+\sum_{j=0}^{k}(2n-2k+j)b^{2n-k}_{k-j}y^{2n-2k+j-1}x^{k-j}\,.$$
	From the equality (\ref{eqi8}), we get 
	$$\phi^{2n-k-1}_x=\frac{P^{2n+1}_{k} (2n-2k-1)}{2}\, ay^{2n-2k-3}x^{k+1}+\sum_{j=0}^{k}\frac{(2n-2k+j)b^{2n-k}_{k-j}}{2}y^{2n-2k+j-2}x^{k-j}\,.$$
	Integrating with respect to $x$, we find a constant $b^{2n-k-1}_{0}\in F$ such that $\phi^{2n-k-1}$ is equal to
	$$P^{2n+1}_{k+1}ay^{2n-2k-3}x^{k+2}+\sum_{j=0}^{k}\frac{(2n-2k+j)b^{2n-k}_{k-j}}{2(k-j+1)}y^{2n-2k+j-2}x^{k-j+1}+b^{2n-k-1}_{0}y^{2n-k-1}\,,$$
	taking into account that $ \displaystyle P^{2n+1}_{k+1}=\frac{P^{2n+1}_{k}\cdot (2n-2k-1)}{2\cdot (k+2)}$. Defining $$\displaystyle b^{2n-k-1}_{k+1-j}=\frac{(2n-2k+j)b^{2n-k}_{k-j}}{2(k-j+1)}$$ for all $0\leq j\leq k $, we get
	$$\phi^{2n-k-1}=P^{2n+1}_{k+1}\, ay^{2n-2k-3}x^{k+2}+\sum_{j=0}^{k+1}b^{2n-k-1}_{k+1-j}y^{2n-2k+j-2}x^{k-j+1}$$
	as wished, and this completes the recurrence. In particular, replacing $k=n-1$ in \eqref{eqodd}, we obtain:
	\begin{equation*}
		\phi^{n+1}=\frac{1\cdot 3\cdots (2n+1)}{2^{n}\cdot n!}ayx^{n}+\sum_{j=0}^{n-1}b^{n+1}_{n-1-j}y^{j+2}x^{n-1-j}\,.
	\end{equation*}
	Considering the equation (\ref{eqi8}) for  $k=n$,  we get $2y\phi^{n}_x=\phi^{n+1}_y$, hence
	\begin{align}%\label{eqq}
		2y\phi^{n}_x&=\frac{1\cdot 3\cdots (2n+1)}{2^{n-1}\cdot n!}ax^{n}+\sum_{j=0}^{n-1}(j+2)b^{n+1}_{n-1-j}y^{j+1}x^{n-1-j}\,.
	\end{align}
	From this equality, we deduce that $a=0$, and consequently $\phi^{2n+1}=0$. Therefore,
	\begin{align*}
		\varphi=z\phi^{2n}+z^2\phi^{2n-1}+\cdots+ z^{2n}\phi^{1}+\phi^0 z^{2n+1}=z\psi\,,
	\end{align*}
	where $\psi\in F[x,y,z]$  is a homogeneous polynomial  of degree $2n$. Finally, arguing as in item ($a$), we deduce that $\psi$ is also a solution.
	
\end{proof}

\end{document}
